\section{Discussion and Limitations}
\label{sec:discussion}

\subsection{Interpretation and Task Validity}

The primary result establishes feasibility at a deliberately narrow boundary. AuthGuard-7702 is a lightweight, decompiler-free prototype that scores delegate runtime bytecode before an EIP-7702 authorization decision. On the primary task---distinguishing USENIX-artifact positives from rule-silent weak negatives---its family-grouped G-DET AUPRC is 0.881. This supports generalization to held-out bytecode-similarity families within the studied corpus. It does not establish universal malicious-delegate detection, and the implemented artifact remains an evaluation-grade scorer rather than a complete wallet defense.

The task-alignment audit is central to interpreting this result. Bare delegation designators are not delegate-runtime inputs, while identical bytecode carrying conflicting labels is incompatible with deterministic bytecode-only classification. We therefore froze the input policy before inspecting rerun outcomes, excluded unresolved designators, and quarantined complete cross-class conflicts without favorable relabeling. Existing family identifiers and fold identities were preserved. Several operating-point metrics changed after this correction, illustrating that input and label compatibility should be established before comparing models. The policy is justified by task correctness, not by the direction of any performance change.

Random splitting gives a complementary warning. AuthGuard's random-minus-family AUPRC gap is approximately 0.094, and the exact-hash blocklist rises from 0.321 to 0.551 under the random diagnostic. The increases for AuthGuard and other clone-sensitive baselines indicate that related bytecode crossing splits can produce optimistic estimates. Family grouping controls this identified leakage channel and yields the more demanding estimate used for primary claims. However, the constructed families are similarity clusters rather than verified attacker groups, and grouping neither removes every possible form of memorization nor supports a general claim about all prior random-split studies.

\subsection{Evasion Findings and Trade-offs}

The stress tests expose both retained signal and important failure modes. Under G-MUT, the sensitive-name approximation falls to zero recall after selector rewriting, whereas AuthGuard retains 0.530 recall at M3. All 727/727 transformed positives passed at each tier as structure-preserving transformations under our opcode-skeleton checker. This is a syntactic repository invariant, not proof of EVM execution or behavioral equivalence. G-VOL is a separate compound test: metadata, address, and selector changes combined with F200 post-\texttt{STOP} flooding reduce AuthGuard recall to 0.130. Those appended bytes are intended to model flooding, but their dynamic unreachability is not formally proved. The result is a serious residual weakness and is not recovered by the distinct G-ADV experiment.

For G-ADV, source-balanced augmentation improves aggregate performance at the tested held-out flooding severity. At pure-M0 F200, fold-mean AUPRC, recall, and FPR change from 0.561, 0.484, and 0.217 to 0.758, 0.727, and 0.174, respectively. The paired family-clustered estimates tell the same aggregate story: pooled recall changes by $+0.253$ (95\% CI $[0.144,0.379]$), FPR by $-0.049$ ($[-0.086,-0.014]$), and AUPRC by $+0.248$ ($[0.177,0.322]$). Opcode-histogram XGBoost-aug has slightly higher F200 recall, while AuthGuard-aug has higher AUPRC and substantially lower FPR, revealing a discrimination and false-positive trade-off rather than dominance on every metric. Augmentation also improves clean and M3 AUPRC and decreases their FPR, but the family-clustered recall intervals for clean M0 and M3 include zero; those recall changes are not statistically resolved. The residual F200 FPR of 0.174 remains material. Collectively, these results support improvement for the specified held-out transformation, not arbitrary adversarial robustness or complete recovery.

\subsection{Evidence and Deployment Boundaries}

Label validity is the principal scientific limitation. Positive labels originate from the USENIX artifact, so the model may learn characteristics associated with that source labeling process. Likewise, \texttt{benign\_cleared} means not flagged by the process, not verified benign ground truth. Exact-bytecode and similarity-family audits reduce leakage but cannot remove this circularity; the primary task remains rule-labeled positives versus rule-silent weak negatives. The sensitive-name approximation accounts for only a small portion of positive labels, and the mutation result indicates that AuthGuard is not merely reproducing that single rule, but neither observation resolves the broader labeling limitation. Independent sourcing left only one truly novel confirmed positive. The corresponding verdict is \textbf{INSUFFICIENT DATA}: the case is anecdotal and cannot support quantitative external generalization.

The evaluated baselines impose another boundary. We did not execute the full USENIX Gigahorse/Datalog pipeline, and do not claim to outperform it. The sensitive-name and external-call methods are local lightweight approximations; the latter obtains high recall by flagging many negatives and is consequently poorly discriminative. A direct comparison with the full analysis pipeline remains future work.

Finally, local timing covers only preloaded-bytecode feature extraction and prediction: mean 3.411\,ms and p95 9.514\,ms on an Apple M1. It excludes wallet parsing, RPC retrieval, caching, warning presentation, and user interaction, and therefore is not end-to-end wallet latency. Operational use would require threshold calibration, reliable code retrieval, model packaging and update procedures, and interface design. Focused next steps are conservative reachability-aware features; larger independently adjudicated malicious and benign EIP-7702 datasets; time-, chain-, and ecosystem-shift evaluation; execution-aware mutation validation; comparison with the full Gigahorse/Datalog pipeline; calibrated selective warnings or escalation to heavier analysis; and wallet-level latency and user-facing evaluation.
